\documentclass[10pt,a4paper]{amsart}
\usepackage[margin=19mm]{geometry}
\usepackage{amsmath,amssymb,mathtools}
\usepackage[scr=rsfs,cal=euler]{mathalfa}
\usepackage{xfrac}
\usepackage[unicode,hidelinks,bookmarks=false]{hyperref}
\newcommand{\ie}{i.e.\ }
\newcommand{\cf}{cf.\ }
\newcommand{\resp}{respectively\ }
\newcommand{\Subsection}{\subsection}
\providecommand{\nobreakdash}{\nobreak\hbox{-}}
\setlength{\emergencystretch}{2em}
\multlinegap0pt
\usepackage{amssymb}
\newtheorem{thm}{Theorem}
\title{Curves, points, incidences and covering}
\author{Arijit Bishnu, Mathew Francis, Pritam Majumder}
\date{Source: arXiv:2507.21758v4 (CC BY 4.0)}
\begin{document}
\maketitle

\noindent\emph{Source and licence.} Curves, points, incidences and covering. Version arXiv:2507.21758v4 (CC BY 4.0), distributed by arXiv under the Creative Commons Attribution 4.0 International licence, http://creativecommons.org/licenses/by/4.0/, as recorded in the arXiv licence metadata for this exact version and shown on the abstract page https://arxiv.org/abs/2507.21758v4. Full licence notice: \texttt{licenses/arxiv-2507.21758-CC-BY-4.0.txt}.\par\smallskip

\noindent\textit{Editorial scope.} Counted unit: the article's thm thm:line-converse (source0.tex lines 231--234). The article's complete original proof of it is delivered with it (source0.tex lines 236--260, one \texttt{proof} environment of 25 lines). No mathematics is written, repaired or re-derived: the statement and the proof are the article's own lines, in the article's own notation, and no retained line is substituted or re-flowed.  No further source-local context is needed: every cross-reference of every retained line resolves inside the two delivered spans.  Nothing was omitted from any retained span, so the package carries no omission record.  No retained line contains a citation, so the wrapper carries no bibliography and no bibliography entry.  No retained line contains a figure environment, an image-inclusion command or drawing code, so no image is delivered and the package declares no asset dependency.  Editorial formatting only, in addition: an AMS \texttt{amsart} preamble that restates the article's own declarations and macros used by the retained lines, namely the environments \texttt{thm} on separate counters, together with the standard packages those lines need.  The retained environments print in this document as Theorem 1 in order of appearance, whereas the article prints the same items with the numbering of its own full document.  The complete native source of the article as distributed in the arXiv e-print for this exact version is bundled unchanged as \texttt{sources/182361/source0.tex}.
\begin{thm}
	\label{thm:line-converse}
	There exists a finite set $P$ of $n^2$ points in $\mathbb{R}^2$ which can be covered with $n$ lines but no subset of $P$ of size $\Omega (n^2)$ can be contained in a projective transformation of a rectangular grid of size $o(n^3)$.
\end{thm}

\begin{proof}		
	Given any two distinct points $p,p'\in\mathbb{R}^2$, we denote by $\ell(p,p')$ the unique line in $\mathbb{R}^2$ that contains both $p$ and $p'$.
	By an $s\times t$ grid, we mean a point set that can be obtained by a projective transform $f$ of the set $[t]\times [s]$. By a ``horizontal line'' of the grid, we mean a line $\ell(f(1,j),f(t,j))$ for some $j\in [s]$, and by a ``vertical line'' of the grid, we mean a line $\ell(f(i,1),f(i,s))$, for some $i\in [t]$. The ``size'' of an $s\times t$ grid is $st$, i.e., the number of points in it. Note that every horizontal line of a grid intersects every vertical line of the grid (since there is a point of the grid that is contained in both of them).
	
	\medskip	
	
	For each $i\in [n]$, let $L_i$ denote the line with equation $y=i$ and let $\mathcal{L}=\{L_i\}_{1\leq i\leq n}$.
	Let $\mathcal{P}$ be the set of points defined as follows. Define $P_1$ to be some set of $n$ distinct points from the line $L_1$. For each $1<i\leq n$, we define $P_i$ to be a set of $n$ distinct points from $L_i$ that do not lie on any of the lines formed by points on other lines, i.e. in $\{\ell(p,p')\colon p\neq p'$ and $p,p'\in\bigcup_{1\leq j\leq i-1} P_j\}$. Let $\mathcal{P}=\bigcup_{1\leq i\leq n} P_i$. Let $m=|\mathcal{P}|$. Note that we have $|\mathcal{L}|=n$. We claim that for any $\mathcal{P}'\subseteq\mathcal{P}$ such that $|\mathcal{P'}|=\Omega(m)=\Omega(n^2)$, any grid that contains all the points of $\mathcal{P}'$ has size $\Omega(n^3)$. 
	
	\medskip		
	
	Note that by our construction, if any line contains two points $p,p'\in\mathcal{P}$ such that $p\in P_i$ and $p'\in P_j$, where $i\neq j$, then $p$ and $p'$ are the only points in $\mathcal{P}$ that are contained in that line. This implies that the following property is satisfied by $\mathcal{P}$ and $\mathcal{L}$.
	\medskip
	
	\noindent (*) Any line in $\mathbb{R}^2$ that contains more than two points in $\mathcal{P}$ belongs to $\mathcal{L}$.
	\medskip
	
	Since every line in $\mathcal{L}$ contains exactly $n$ points of $\mathcal{P}$, we then have another property.
	\medskip
	
	\noindent (+) Any line in $\mathbb{R}^2$ contains at most $n$ points in $\mathcal{P}$.
	\medskip
	
	Let $\mathcal{P}'\subseteq\mathcal{P}$ be such that $|\mathcal{P}'|=\Omega(n^2)$. Consider any grid $\mathbb{G}$ that contains all the points of $\mathcal{P}'$. Let $\mathbb{G}$ be an $s\times t$ grid. Let $h_1,h_2,\ldots,h_s$ denote the horizontal lines of $\mathbb{G}$ and let $v_1,v_2,\ldots,v_t$ denote the vertical lines of $\mathbb{G}$. Suppose for the sake of contradiction that there exist $i\in [s]$ and $j\in [t]$ such that both the lines $h_i$ and $v_j$ contain at least 3 points of $\mathcal{P}$ each. Then by property~(*), $h_i$ and $v_j$ are both lines in $\mathcal{L}$. But as $h_i$ and $v_j$ intersect, they are two lines in $\mathcal{L}$ that intersect, which is a contradiction, since the lines in $\mathcal{L}$ are all parallel to each other (Note that parallel lines under a projective transformation may not be parallel but they do not intersect at any of the $s\times t$ grid points defined.). Thus, we can conclude without loss of generality that for each $i\in [s]$, the horizontal line $h_i$ of $\mathbb{G}$ contains at most two points from $\mathcal{P}$, and hence at most two points from $\mathcal{P}'$. Since every point in $\mathcal{P}'$ is contained in at least one horizontal line of $\mathbb{G}$, we have that $s\geq |\mathcal{P}'|/2$ and therefore $s=\Omega(n^2)$. By property~(+), each vertical line of $\mathbb{G}$ can contain at most $n$ points of $\mathcal{P}'$, and therefore, $t\geq |\mathcal{P}'|/n$, which implies that $t=\Omega(n)$. Thus the size of the grid $\mathbb{G}$ is $st=\Omega(n^3)$.
\end{proof}

\end{document}
